\section{문제와 동기}
\label{sec:problem}

\paragraph{한 모델이 모든 짐을 진다.} 기존 VLA는 안의 VLM(3--8B)이 과제 계획, 진행 판단, 다음 명령, 실제 움직임을 전부 맡는다. 이 크기의 VLM에게는 너무 무거운 짐이다. 게다가 판단은 VLM의 언어·시각 표현에서 이루어지고 움직임은 행동 전문가가 만들기 때문에, 그 사이의 전이에서 정보가 손실된다.

\paragraph{큰 VLM은 판단은 잘하지만 느리다.} 판단을 큰 VLM에게 맡기면 정확도는 오르지만 한 번 답하는 데 초 단위가 걸린다. 큰 VLM 혼자 제어하면 명령과 명령 사이에 로봇이 새 정보를 보지 못하는 `장님 구간'이 생기고, 1단계 구현에서는 로봇이 답을 기다리며 멈춘다(\cref{fig:timeline}a).

\paragraph{목표.} 큰 VLM의 판단력과 작은 JCR의 빠른 반응을 \emph{성능 손실 없이} 잇는다. 이를 위해 (i) 계획·판단·명령은 상위 VLM에 모으고, (ii) JCR에게는 상위 명령 기준의 미세 조종이라는 단순한 일만 주며, (iii) 두 모듈을 잇는 과정에서 새는 성능을 직접 재고 줄인다. 따라서 JCR의 설계·학습·평가는 JCR가 스스로 계획하는 쪽으로 가지 않는다. JCR 입력의 목표는 상위 명령이고, 출력은 허용 범위 안의 보정과 이상 신호뿐이다.
